\section{Conclusion}
\label{sec:conclusion}

We benchmarked six neural architectures that share a 385-parameter budget and differ only in how those parameters
are arranged. Three findings stand out. First, configuration changes final accuracy only when the target has
structure for it to match. A staged modular hybrid reduced error on a compositional task by \chainHBGain\%, a gap that
held after 300 epochs, while every other structure tied on simple tasks. Second, small-model latency is decided by
operator count at small batch sizes and by activation width at large ones, and parameter count predicts neither;
\texttt{torch.compile} narrows the operator-count penalty but adds a fixed per-call cost. Third, protocol choices can
masquerade as architecture effects: a cross-framework confound in our pilot and unscaled targets in an earlier run
both produced large apparent advantages that disappeared under a controlled protocol.

Future work will extend the study to other runtimes (ONNX Runtime) and hardware (GPUs, neural accelerators,
microcontrollers), to larger parameter budgets where arithmetic dominates overhead, and to real-world tasks where
the stage decomposition must be learned rather than designed.
